\documentclass{article}
\usepackage{forest,ctex,amsmath,color,xcolor,colortbl,graphicx,geometry,comment,listings,hyperref,booktabs,tabularx,wrapfig,caption,calc,mhchem,fix-cm,tocloft,titlesec,setspace,fancyhdr,booktabs,float,subcaption,enumitem,pdfpages}
\usepackage{lastpage}
\usepackage[utf8]{inputenc}
\usepackage[absolute,overlay]{textpos}  % 绝对定位
\setlength{\parindent}{2em} 

\geometry{left=2.5cm,right=2.5cm,top=2.5cm,bottom=2.5cm}
%\setlength{\parindent}{2pt} 
\date{}
\title{\textbf{AI 工具使用详情}}
\begin{document}
\maketitle
\section{所用 AI 工具名称和版本}
1.Deepseek V3

2.豆包 2024.05 核心能力基线版

\section{具体使用目的和环节}
1.搜索算法模型，并自行判断其对于带解决问题的适配性

2.对于不熟悉库与代码进行询问以便利用其进行算法编写

\section{关键交互记录（重要提示词与回复）}
\subsection{问题一}
Q：如何解读 GAMM 模型中 “Y 浓度与 BMI、孕周关系” 的输出结果？

A：先看模型基础信息，如 “Family: gaussian” 表示 Y 浓度服从正态分布，“Link function:
identity” 说明直接建模均值；参数系数中，截距项显著（p<2e-16）代表存在 Y 浓度的基础水平；
平滑项里，s (BMI) 和 s (Gestational\_days) 的 p 值 < 0.05，说明二者对 Y 浓度的非线性效应
显著，edf 值越大非线性越强；拟合优度看调整后 R²（如 0.753）和偏差解释率（如 81\%），数值越高
模型拟合效果越好。

Q：用 GAMM 建模时，如何选择 “连接函数”？

A：连接函数的选择取决于因变量（响应变量）的类型：若因变量是连续型（如 Y 染色体浓度、血压），选
“identity”（恒等连接），直接建模均值；若因变量是二分类（如是否患病、是否购买），选 “logit”
（逻辑连接），建模概率的对数优势比；若因变量是计数型（如发病次数、点击次数），选 “log”（对数
连接），建模速率的对数。

Q：用 GAMM 分析二元关系时，如何判断“非线性效应是否显著”？

A：主要看 GAMM 输出结果中 “平滑项的近似显著性” 部分：对于 “s(x)”，首先看 p 值，若 p<0.05
（如 p<2e-16），说明平滑效应（非线性效应）显著；其次看有效自由度（edf），若 edf>1，说明二元
的关系是非线性的（edf 值越大，非线性程度越强），若 edf≈1，则接近线性关系，可考虑用线性模型替代。

Q：GAMM 模型的 “拟合优度指标” 有哪些？分别代表什么含义？

A：主要有两个核心指标：一是 “调整后 R²（R-sq.(adj)）”，它衡量模型能解释因变量变异的比例，
如 0.753 表示模型解释了 Y 浓度 75.3\% 的变异，该值越高拟合越好；二是 “偏差解释率（Deviance
explained）”，偏差是衡量模型预测值与真实值差异的指标，偏差解释率越高，说明模型对数据偏差的解
释能力越强，如 81\% 表示模型解释了 81\% 的偏差。

\subsection{问题二}
Q：肘部法则是什么？它的核心作用是什么？

A：肘部法则是聚类分析（如 K-Means）中确定最优聚类数的启发式方法。核心是计算不同聚类数 K 对应的组内平方和（WCSS），绘制 “K-WCSS” 曲线，找到 WCSS 下降速率由快转慢的 “肘部”（拐点），该拐点对应的 K 即为最优聚类数，可平衡聚类效果与计算成本，避免过度聚类或聚类不足。

Q：用肘部法则确定 K 值的具体步骤是什么？

A：第一步，对聚类变量做标准化（消除量纲影响）；第二步，遍历 K=1 到 K=10（或根据业务设定范围），每个 K 训练 K-Means 模型并记录 WCSS；第三步，绘制 “K-WCSS” 曲线；第四步，找到曲线中 WCSS 下降由快转慢的拐点，该拐点对应的 K 就是最优聚类数，如 K=3 时曲线从快速下降转为平缓，即选 K=3。
\subsection{问题三}
Q：用肘部法则确定 K 值时，为什么要对 “聚类变量做标准化”？

A：因为聚类算法（如 K-Means）是基于 “距离” 计算的，若变量的量纲不同，量纲大的变量会主导距离计算，
影响聚类结果。例如，“身高（cm，数值范围 150-190）” 和 “体重（kg，数值范围 40-100）”，若不标准化，
身高的数值差异（40）远小于体重的数值差异（60），但实际二者对聚类的重要性可能相同，标准化后（将
变量转为均值 0、标准差 1），可消除量纲影响，让各变量在聚类中发挥同等作用。
\subsection{问题四}
Q：数学建模中的 “shaq” 可能指什么？

A：结合建模场景，“shaq” 更可能是拼写偏差，对应 SHA（安全哈希算法）。SHA 是密码学中的哈希函数，
能将任意长度数据转为固定长度哈希值，在建模中可用于敏感数据脱敏（如用哈希值替换患者身份证号）、
数据完整性校验（比对哈希值判断数据是否被篡改）、高维数据特征压缩（将文本转为哈希值降维），常见
于医疗、金融等需数据安全的建模场景。

Q：随机森林中 “Bootstrap 抽样” 的具体过程是什么？

A：Bootstrap 抽样是有放回的抽样方式，具体过程为：假设原始训练集有 N 个样本，每次从 N 个样本中随机抽取 1 个样本，记录后将其放回原始样本集，重复该过程 N 次，得到一个包含 N 个样本的子样本集（部分样本可能被多次抽取，部分样本可能未被抽取）。每棵决策树都基于一个通过 Bootstrap 抽样得到的子样本集训练，这保证了树与树之间的差异性。

Q：随机森林适合哪些建模场景？

A：适合分类和回归两类场景。分类场景如客户流失预测（判断客户是否会流失）、疾病诊断（根据体检指
标判断是否患病）；回归场景如房价预测（结合面积、地段等预测房价）、销量预测（根据广告投入、季节
等预测销量）。此外，还可用于特征重要性分析，筛选对结果影响最大的变量，辅助模型优化。

Q：点二列相关系数与 “二列相关系数” 的区别是什么？

A：核心区别在于二分变量的类型：点二列相关系数适用于 “真正二分变量”，即变量本身就是二分的（如性
别、是否患病），不是由连续变量人为划分的；二列相关系数适用于 “人为二分变量”，即变量原本是连续的，
后被人为划分为二分（如将 “成绩” 划分为 “及格 / 不及格”，将 “收入” 划分为 “高收入 / 低收入”）。
此外，点二列相关系数的计算公式更简单，二列相关系数则需用到连续变量的正态分布假设。

\section{采纳和人工修改情况}

\end{document}